% Options for packages loaded elsewhere
\PassOptionsToPackage{unicode}{hyperref}
\PassOptionsToPackage{hyphens}{url}
\documentclass[
  a4paper,11pt]{article}
\usepackage{xcolor}
\usepackage[margin=1in]{geometry}
\usepackage{amsmath,amssymb}
\setcounter{secnumdepth}{-\maxdimen} % remove section numbering
\usepackage{iftex}
\ifPDFTeX
  \usepackage[T1]{fontenc}
  \usepackage[utf8]{inputenc}
  \usepackage{textcomp} % provide euro and other symbols
\else % if luatex or xetex
  \usepackage{unicode-math} % this also loads fontspec
  \defaultfontfeatures{Scale=MatchLowercase}
  \defaultfontfeatures[\rmfamily]{Ligatures=TeX,Scale=1}
\fi
\usepackage{lmodern}
\ifPDFTeX\else
  % xetex/luatex font selection
\fi
% Use upquote if available, for straight quotes in verbatim environments
\IfFileExists{upquote.sty}{\usepackage{upquote}}{}
\IfFileExists{microtype.sty}{% use microtype if available
  \usepackage[]{microtype}
  \UseMicrotypeSet[protrusion]{basicmath} % disable protrusion for tt fonts
}{}
\makeatletter
\@ifundefined{KOMAClassName}{% if non-KOMA class
  \IfFileExists{parskip.sty}{%
    \usepackage{parskip}
  }{% else
    \setlength{\parindent}{0pt}
    \setlength{\parskip}{6pt plus 2pt minus 1pt}}
}{% if KOMA class
  \KOMAoptions{parskip=half}}
\makeatother
\usepackage{longtable,booktabs,array}
\newcounter{none} % for unnumbered tables
\usepackage{calc} % for calculating minipage widths
% Correct order of tables after \paragraph or \subparagraph
\usepackage{etoolbox}
\makeatletter
\patchcmd\longtable{\par}{\if@noskipsec\mbox{}\fi\par}{}{}
\makeatother
% Allow footnotes in longtable head/foot
\IfFileExists{footnotehyper.sty}{\usepackage{footnotehyper}}{\usepackage{footnote}}
\makesavenoteenv{longtable}
\setlength{\emergencystretch}{3em} % prevent overfull lines
\providecommand{\tightlist}{%
  \setlength{\itemsep}{0pt}\setlength{\parskip}{0pt}}
\usepackage{bookmark}
\IfFileExists{xurl.sty}{\usepackage{xurl}}{} % add URL line breaks if available
\urlstyle{same}
\hypersetup{
  hidelinks,
  pdfcreator={LaTeX via pandoc}}

\author{}
\date{}

\begin{document}

\section{Geometric Structure of Generation Mixing --- Qualitative
Derivation of the CKM and PMNS
Matrices}\label{geometric-structure-of-generation-mixing-qualitative-derivation-of-the-ckm-and-pmns-matrices}

\subsection{------ A Structural Answer to the Generation Mixing Problem
of {[}4{]} §9.11.6
------}\label{a-structural-answer-to-the-generation-mixing-problem-of-4-9.11.6}

\textbf{Author}: Noriaki Kihara (WF System Co., Ltd.)

\textbf{ORCID}: https://orcid.org/0009-0002-7952-9855

\textbf{DOI}: https://doi.org/10.5281/zenodo.19731602

\textbf{Date}: April 2026

\begin{center}\rule{0.5\linewidth}{0.5pt}\end{center}

\subsection{§0 Purpose of This Paper}\label{purpose-of-this-paper}

In {[}4{]} §9.11.6, we declared that ``how generation mixing (the CKM
matrix and the PMNS matrix) is derived as the breaking of \(S_3\)
symmetry of the \(xyz\) axes remains an open problem.''

The present paper responds to this declaration and demonstrates the
following.

\begin{enumerate}
\def\labelenumi{\arabic{enumi}.}
\tightlist
\item
  The diagonally dominant structure of the CKM matrix follows directly
  from the \(\mathrm{SO}(3) \to \mathrm{SO}(2)\) symmetry breaking
  caused by the Higgs (standing wave) axis selection (§1).
\item
  The large mixing angle structure of the PMNS matrix follows directly
  from the (spatial, spatial)-type 2-face structure of neutrinos (§2).
\item
  The structural difference between the CKM and PMNS matrices reduces to
  the geometric difference between (spatial, \(t\))-type and (spatial,
  spatial)-type (§3).
\end{enumerate}

The claims of this paper are limited to the derivation of the
\textbf{qualitative structure} of the mixing matrices (diagonal
dominance vs.~large mixing angles). The derivation of specific numerical
values of the mixing angles is an open problem sharing the same root as
the quantitative derivation of mass ratios ({[}4{]} §9.11.4) and lies
outside the scope of this paper.

\textbf{References}: - {[}4{]} Set Structure and Combinatorial
Properties of the 6-Dimensional Hyperrectangle (W8) - {[}5{]}
Interactions of Shape-Invariant Waves (W11)

\begin{center}\rule{0.5\linewidth}{0.5pt}\end{center}

\subsection{\texorpdfstring{§1 The CKM Matrix: Generation Mixing of the
(Spatial,
\(t\))-Type}{§1 The CKM Matrix: Generation Mixing of the (Spatial, t)-Type}}\label{the-ckm-matrix-generation-mixing-of-the-spatial-t-type}

\subsubsection{1.1 Correspondence Between Generations and Spatial
Axes}\label{correspondence-between-generations-and-spatial-axes}

By {[}4{]} §9.1, charged fermions (quarks and charged leptons) have
2-faces of the (spatial 1-axis, \(t\))-type. The three spatial axes
\(x, y, z\) correspond to the three generations:

{\def\LTcaptype{none} % do not increment counter
\begin{longtable}[]{@{}cccc@{}}
\toprule\noalign{}
Generation & Spatial axis & Up-type quark & Down-type quark \\
\midrule\noalign{}
\endhead
\bottomrule\noalign{}
\endlastfoot
1 & \(x\) & \(u = (+, 0, 0, +, 0, Q)\) & \(d = (+, 0, 0, -, 0, Q)\) \\
2 & \(y\) & \(c = (0, +, 0, +, 0, Q)\) & \(s = (0, +, 0, -, 0, Q)\) \\
3 & \(z\) & \(t = (0, 0, +, +, 0, Q)\) & \(b = (0, 0, +, -, 0, Q)\) \\
\end{longtable}
}

\subsubsection{1.2 Axis Selection by the Higgs and Symmetry
Breaking}\label{axis-selection-by-the-higgs-and-symmetry-breaking}

By {[}4{]} §9.11.7, the Higgs boson (standing wave, \(\kappa = 1\))
takes a nonzero component along one of the three \(xyz\) axes. In the
initial state, the \(xyz\) axes are completely equivalent under
\(\mathrm{SO}(3)\) symmetry, but the symmetry is spontaneously broken
when the Higgs selects one axis. We denote the selected axis as \(z\).

As a result: - \(z\) axis: the nonzero axis of the Higgs → direct
coupling → largest mass (third generation) - \(x, y\) axes: do not share
the nonzero axis of the Higgs → only indirect coupling → lighter masses
(first and second generations)

The symmetry is broken as \(\mathrm{SO}(3) \to \mathrm{SO}(2)\), where
the \(z\) axis is separated while an approximate \(\mathrm{SO}(2)\)
symmetry remains between the \(x\) and \(y\) axes.

\subsubsection{1.3 Diagonally Dominant Structure of the CKM
Matrix}\label{diagonally-dominant-structure-of-the-ckm-matrix}

\textbf{Proposition 1.1}. In generation mixing of (spatial, \(t\))-type
fermions, the Higgs \(z\)-axis selection leads to the following
hierarchical structure:

\begin{enumerate}
\def\labelenumi{(\roman{enumi})}
\item
  Mixing between the third generation and the first/second generations
  is small (separation of the \(z\) axis).
\item
  Mixing between the first and second generations is larger than (i)
  (approximate symmetry of the \(x, y\) axes).
\item
  The CKM matrix is diagonally dominant.
\end{enumerate}

\emph{Argument}. The weak interaction mediated by \(W^{\pm}\) is an
isospin transition ({[}5{]} §6), which acts on all position-type axes
via the axis-type symmetry, flipping the \(t\)-axis sign of the fermion
and converting up-type quarks into down-type quarks. Since each
generation of quarks occupies only one spatial axis, inter-generation
transitions involve changing the spatial axis.

Due to the Higgs \(z\)-axis selection, fermions on the \(z\) axis (third
generation) have a mass scale significantly different from those on
other axes. The deviation between mass eigenstates and flavor
eigenstates (mixing angle) depends on the coupling strength between
axes, which is inversely proportional to the mass difference.

\begin{enumerate}
\def\labelenumi{(\roman{enumi})}
\item
  Since \(m_3 \gg m_{1,2}\), the mixing angles between the third
  generation and the first/second generations are small.
\item
  Since the \(x\) and \(y\) axes are approximately equivalent under
  \(\mathrm{SO}(2)\) symmetry, \(m_1 \approx m_2\), and the mixing angle
  between the first and second generations is larger than in (i).
\item
  From (i) and (ii), the diagonal elements of the CKM matrix are close
  to 1 and the off-diagonal elements are small (diagonal dominance).
  \(\square\)
\end{enumerate}

\textbf{Comparison with experimental values}:

\[V_{\text{CKM}} \approx \begin{pmatrix} 0.974 & 0.225 & 0.004 \\ 0.225 & 0.973 & 0.041 \\ 0.009 & 0.040 & 0.999 \end{pmatrix}\]

\begin{itemize}
\tightlist
\item
  \(|V_{ub}| = 0.004\), \(|V_{td}| = 0.009\): mixing between the 3rd and
  1st generations is smallest ✓
\item
  \(|V_{cb}| = 0.041\), \(|V_{ts}| = 0.040\): mixing between the 3rd and
  2nd generations is small ✓
\item
  \(|V_{us}| = 0.225\): mixing between the 1st and 2nd generations is
  largest ✓
\end{itemize}

\begin{center}\rule{0.5\linewidth}{0.5pt}\end{center}

\subsection{§2 The PMNS Matrix: Generation Mixing of the (Spatial,
Spatial)-Type}\label{the-pmns-matrix-generation-mixing-of-the-spatial-spatial-type}

\subsubsection{2.1 The 2-Face Structure of
Neutrinos}\label{the-2-face-structure-of-neutrinos}

By {[}4{]} §9.9 Table A, neutrinos have 2-faces of the (spatial,
spatial)-type. The \(\binom{3}{2} = 3\) combinations of spatial axes
correspond to the three flavors:

{\def\LTcaptype{none} % do not increment counter
\begin{longtable}[]{@{}ccl@{}}
\toprule\noalign{}
Flavor & Nonzero axes & Set \\
\midrule\noalign{}
\endhead
\bottomrule\noalign{}
\endlastfoot
\(\nu_e\) & \((x, y)\) & \((+, +, 0, 0, 0, 0)\) \\
\(\nu_\mu\) & \((y, z)\) & \((0, +, +, 0, 0, 0)\) \\
\(\nu_\tau\) & \((x, z)\) & \((+, 0, +, 0, 0, 0)\) \\
\end{longtable}
}

\subsubsection{2.2 Axis-Sharing Structure}\label{axis-sharing-structure}

\textbf{Proposition 2.1}. Any two neutrino flavors share exactly one
spatial axis.

\emph{Proof}. Among the three combinations of choosing 2 elements from
the 3-element set \(\{x, y, z\}\) --- namely
\(\{x,y\}, \{y,z\}, \{x,z\}\) --- the intersection of any two
combinations consists of exactly one element:

{\def\LTcaptype{none} % do not increment counter
\begin{longtable}[]{@{}ccc@{}}
\toprule\noalign{}
Pair & Shared axis & Non-shared axes \\
\midrule\noalign{}
\endhead
\bottomrule\noalign{}
\endlastfoot
\(\nu_e, \nu_\mu\) & \(y\) & \(x \leftrightarrow z\) \\
\(\nu_\mu, \nu_\tau\) & \(z\) & \(y \leftrightarrow x\) \\
\(\nu_e, \nu_\tau\) & \(x\) & \(y \leftrightarrow z\) \\
\end{longtable}
}

This is a combinatorial property of \(\binom{3}{2}\) and holds in
general for the case of 3 axes with 2 chosen. \(\square\)

\subsubsection{2.3 Large Mixing Angle Structure of the PMNS
Matrix}\label{large-mixing-angle-structure-of-the-pmns-matrix}

\textbf{Proposition 2.2}. In generation mixing of (spatial,
spatial)-type fermions, the Higgs \(z\)-axis selection leads to the
following structure:

\begin{enumerate}
\def\labelenumi{(\roman{enumi})}
\item
  Both \(\nu_\mu\) and \(\nu_\tau\) contain the \(z\) axis (Higgs axis),
  and their non-shared axes \((y, x)\) are unaffected by the Higgs, so
  \(\nu_\mu \leftrightarrow \nu_\tau\) mixing is nearly maximal.
\item
  Since \(\nu_e\) does not contain the \(z\) axis, the mixing
  \(\nu_e \leftrightarrow \nu_\mu\) and
  \(\nu_e \leftrightarrow \nu_\tau\) is smaller than (i).
\item
  The PMNS matrix has large off-diagonal elements (in contrast to the
  CKM matrix).
\end{enumerate}

\emph{Argument}.

\begin{enumerate}
\def\labelenumi{(\roman{enumi})}
\item
  \(\nu_\mu = (y, z)\) and \(\nu_\tau = (x, z)\) share the axis \(z\),
  with non-shared axes \(y\) and \(x\), respectively. Since \(x\) and
  \(y\) are unaffected by the Higgs and the \(\mathrm{SO}(2)\) symmetry
  remains, the coupling strengths of \(\nu_\mu\) and \(\nu_\tau\) to the
  Higgs are nearly equal. The mixing angle of two states with equal
  coupling strengths approaches the maximum (\(\theta \approx 45°\)).
\item
  \(\nu_e = (x, y)\) does not contain the \(z\) axis. The coupling of
  \(\nu_e\) to the Higgs is limited to indirect coupling through the
  \(x, y\) axes, and its coupling structure differs from that of
  \(\nu_\mu, \nu_\tau\), which contain the \(z\) axis. This asymmetry
  causes the mixing angle of \(\nu_e\) with \(\nu_{\mu,\tau}\) to
  deviate from the maximum.
\item
  The coexistence of the maximal mixing of (i) and the moderate mixing
  of (ii) results in large off-diagonal elements in the PMNS matrix
  overall. \(\square\)
\end{enumerate}

\textbf{Comparison with experimental values}:

\begin{itemize}
\tightlist
\item
  \(\theta_{23} \approx 49°\) (atmospheric angle,
  \(\nu_\mu \leftrightarrow \nu_\tau\)): nearly maximal ✓ ---
  Proposition (i)
\item
  \(\theta_{12} \approx 33°\) (solar angle,
  \(\nu_e \leftrightarrow \nu_\mu\)): large but not maximal ✓ ---
  Proposition (ii)
\item
  \(\theta_{13} \approx 8.6°\) (reactor angle,
  \(\nu_e \leftrightarrow \nu_\tau\)): smallest ✓ --- Proposition (ii)
\end{itemize}

\begin{center}\rule{0.5\linewidth}{0.5pt}\end{center}

\subsection{§3 Origin of the Structural
Difference}\label{origin-of-the-structural-difference}

\subsubsection{\texorpdfstring{3.1 (Spatial, \(t\))-Type vs.~(Spatial,
Spatial)-Type}{3.1 (Spatial, t)-Type vs.~(Spatial, Spatial)-Type}}\label{spatial-t-type-vs.-spatial-spatial-type}

\textbf{Proposition 3.1}. The structural difference between the CKM
matrix (diagonally dominant) and the PMNS matrix (large mixing angles)
reduces to the difference in axis composition of the 2-faces.

{\def\LTcaptype{none} % do not increment counter
\begin{longtable}[]{@{}
  >{\raggedright\arraybackslash}p{(\linewidth - 4\tabcolsep) * \real{0.3333}}
  >{\raggedright\arraybackslash}p{(\linewidth - 4\tabcolsep) * \real{0.3333}}
  >{\raggedright\arraybackslash}p{(\linewidth - 4\tabcolsep) * \real{0.3333}}@{}}
\toprule\noalign{}
\begin{minipage}[b]{\linewidth}\raggedright
Property
\end{minipage} & \begin{minipage}[b]{\linewidth}\raggedright
(Spatial, \(t\))-type
\end{minipage} & \begin{minipage}[b]{\linewidth}\raggedright
(Spatial, spatial)-type
\end{minipage} \\
\midrule\noalign{}
\endhead
\bottomrule\noalign{}
\endlastfoot
Number of occupied spatial axes & \textbf{1} & \textbf{2} \\
Generations containing the Higgs axis (\(z\)) & Third generation only &
\textbf{Two}: \(\nu_\mu\) and \(\nu_\tau\) \\
Axis sharing between two flavors & None (each generation has an
independent axis) & Always share 1 axis (Proposition 2.1) \\
Separation by the Higgs & Split into 1 vs.~2 (\(z\) vs.~\(x, y\)) & Weak
separation (2 containing \(z\) vs.~1 not containing \(z\)) \\
Consequence for mixing angles & Small (diagonally dominant) & Large \\
\end{longtable}
}

\emph{Argument}. In the (spatial, \(t\))-type, each generation occupies
one independent spatial axis, with no axis sharing between generations.
The Higgs \(z\)-axis selection separates only one generation, and the
coupling between the remaining two generations is also indirect, so the
mixing angles are small overall.

In the (spatial, spatial)-type, each flavor occupies two spatial axes,
and any two flavors share one axis (Proposition 2.1). The Higgs \(z\)
axis is contained in two out of three flavors, so the separation effect
is weak, and the direct coupling through shared axes makes the mixing
angles large. \(\square\)

\subsubsection{3.2 Relation to Supplementary Lecture
1}\label{relation-to-supplementary-lecture-1}

As shown in Supplementary Lecture 1 §2.2, (spatial, spatial)-type
neutrinos have \(\text{sign}(M_F) = +1\) (positive orientation), which
differs in orientation from (spatial, \(t\))-type charged fermions. By
Supplementary Lecture 1 Proposition 2.2, antineutrinos also have
\(M_F^{\text{anti}} = M_F\) (sign-invariant).

The structural difference in generation mixing shares the same geometric
origin as this difference in orientation --- the difference in 2-face
structure between the (spatial, \(t\))-type and the (spatial,
spatial)-type.

\begin{center}\rule{0.5\linewidth}{0.5pt}\end{center}

\subsection{§4 Conclusions}\label{conclusions}

In this paper, we have provided the following structural answers to the
generation mixing problem declared open in {[}4{]} §9.11.6.

\begin{enumerate}
\def\labelenumi{\arabic{enumi}.}
\item
  \textbf{The diagonally dominant structure of the CKM matrix} follows
  from the fact that each generation of (spatial, \(t\))-type fermions
  occupies an independent spatial axis, and the Higgs \(z\)-axis
  selection induces \(\mathrm{SO}(3) \to \mathrm{SO}(2)\) symmetry
  breaking that separates the third generation (Proposition 1.1).
\item
  \textbf{The large mixing angle structure of the PMNS matrix} follows
  from the fact that each flavor of (spatial, spatial)-type neutrinos
  occupies two spatial axes, and any two flavors share one axis
  (Proposition 2.1). In particular, since both \(\nu_\mu\) and
  \(\nu_\tau\) contain the Higgs axis \(z\), \(\theta_{23}\) is nearly
  maximal (Proposition 2.2).
\item
  \textbf{The structural difference between the CKM and PMNS matrices}
  reduces to the difference in 2-face structure between the (spatial,
  \(t\))-type and the (spatial, spatial)-type (Proposition 3.1).
\end{enumerate}

The derivation of specific numerical values of the mixing angles remains
as an open problem sharing the same root as the quantitative derivation
of mass ratios in {[}4{]} §9.11.4.

\begin{center}\rule{0.5\linewidth}{0.5pt}\end{center}

\subsection{References}\label{references}

\begin{itemize}
\tightlist
\item
  {[}4{]} Noriaki Kihara, ``Set Structure and Combinatorial Properties
  of the 6-Dimensional Hyperrectangle,'' Zenodo preprint, DOI:
  10.5281/zenodo.19762134 (2026).
\item
  {[}5{]} Noriaki Kihara, ``Interactions of Shape-Invariant Waves ---
  Axial Displacement Transfer, Wave Packet Deformation, and Retroactive
  Constitution of Causality,'' Zenodo preprint, DOI:
  10.5281/zenodo.19763463 (2026).
\end{itemize}

\end{document}
